\section{Chapter Summary}

This chapter detailed the complete technical implementation of the DBaaS platform, bridging the gap between high-level architectural design and the underlying cloud-native reality. By utilizing Go for its exceptional concurrency and native Kubernetes integration, the backend successfully enforces strict fault isolation between the Control Plane (REST API) and the Data Plane (PGProxy). Furthermore, delegating database lifecycle management to declarative Kubernetes operators allows the system to achieve highly efficient resource utilization, scaling seamlessly across cloud environments without vendor lock-in.

The integration of Flutter as the presentation layer ensures that this complex orchestration is safely abstracted behind a responsive, unidirectional data flow, empowering users to intuitively provision, manage, and query relational and NoSQL databases. With the infrastructure containerized via ultra-lightweight Alpine Docker images and deployed through an automated CI/CD pipeline, the platform stands as a highly resilient, performant, and production-ready solution capable of handling multi-tenant database operations at scale.
